\section{Proofs: IOPP soundness \textcolor{red}{(tbc)}}

\subsection{Protocol \ref{prot:main}}
Let us outline the soundness properties of the protocol.
In the following we write
\begin{equation}
\ell_1 = \ell_1(\theta_1)
\end{equation}
for the Johnson list size bound on the number of proximates from $C_1$ at distance $\theta_1$.
% For both $f_1$ and $f_2$, we define $S(f_1)$ and $S(f_2)$ to be the set of $z\in F$, which distinguish all $\theta_1$- and $\theta_2$ proximates by value. 
% Note that the exceptional set 
% \begin{equation}
% \left|F\setminus S(f_i)\right| \leq \binom{\ell_i}{2} \cdot 2^{m_i},
% \end{equation} 
% %contains at most $\ell_{\theta_i}^2 \cdot 2^{m_1}$ many points, 
% the union bound of common zeros of any pair of proximates.
% We define
% \begin{equation*}
% \mathscr L_{k}\big|_{z,v} = \left\{p(X)\in \mathscr L_k \::\: p(z) = v\right\}.
% \end{equation*}
% % \begin{equation*}
% % E_z: \mathscr L_{m_1} \cap B_\theta(f_1) \longrightarrow F, \quad p(X) \mapsto p(z)
% % \end{equation*}

\paragraph{Phase \ref{i:main:phase:1}:}
The first phase is a reduction $\mathcal R_\theta \longrightarrow \mathcal R_\text{phase 1}$ depending on the verifier challenges $\vec\xi\sample \mathbb F_q^{t}$, where
\begin{align*}
\mathcal R_\text{phase 1} 
&=
\left\{
    \begin{matrix}
        f_1 \in \left(\mathbb F_q^{D}\right)^{e},% v_1\in F 
        \\
        f_2^* \in \left(\mathbb F_q^{D}\right)^{e}%, v_2\in F 
    \end{matrix}
    \::\:
    % \hspace*{-0.25cm}
    \begin{matrix}
        \exists (p_1(X),p_2(X)) \in \mathscr P_{m_1}^2 %\times \mathscr P_{m_1}
        \\
        d((p_1,p_2), (f_1,f_2^*)) \leq \theta_1,
        \\
        \text{and}
        \\
        (p_{1}(X), \vec\xi, p_{2}(X))
        \\
        \text{satisfy the lifted R1CS} 
        \\
        \text{on input }\vec{in}
        % \text{and their coefficient vectors}
        % \\
        % (\vec c_1, \vec c_2) \text{ satisfy the lifted R1CS }
        % \\
        % \text{on inputs }\vec{in} , \vec\xi 
    \end{matrix}
\right\},
\end{align*}
where as before, the distance refers to the interleaved representations of the polynomials.
The phase is a single-step reduction, and its soundness error is as follows.

\begin{lem}
Phase \ref{i:main:phase:1} in Protocol \ref{prot:main} is a round-by-round sound oracle reduction $\mathscr R_{\theta_1}\longrightarrow R_\text{phase 1}$ with soundness error
\[
\varepsilon_\text{phase 1} \leq 
\ell_1\cdot \varepsilon_\text{R1CS}, 
\]
where $\varepsilon_\text{R1CS}$ is the soundness error of the randomized R1CS.
\end{lem}

\begin{proof}
Assume that for more than $\kappa = \ell_1(\theta )\varepsilon_\text{R1CS}\cdot |\mathbb F_q|^t$ distinct verifier challenges $\vec\xi$, there exists a potential prover answer $f_2\in \mathbb F_q^{D'}$ which extends $f_1$ to a pair $(f_1,f_2^*)$ in $\mathscr R_\text{phase 1}$.
For each such challenge $\vec\xi_i$, $i=1,\ldots, \kappa$, we select one of the $\theta_1$-proximates $(p_{1}(X), p_{2}(X))$ stated by $\mathscr R_\text{phase 1}$, and obtain a list 
%$\mathscr L$ of more than $\kappa$ many $(p_{1,i}(X),\vec\xi_i, p_{2,i}(X))$.
\[
\mathscr L = \left\{
(p_{1,i}(X),\vec\xi_i , p_{2,i}(X))) \::\: 
\begin{matrix}
d(p_{1,i}(X),f_1)\leq \theta_1,
\\\text{and}
\\
(p_{1,i}(X), \vec\xi_i, p_{2,i}(X))
\\
\text{satisfy the lifted R1CS} 
\\
\text{on input }\vec{in}
\end{matrix}
\right\}
\]
with more than $\kappa$ entries. 

By the list size bound, there at most $\ell_1$ distinct $p_{1,i}(X)$ mentioned in $\mathscr L$.
One of them, $p_1(X) = p_{1,i_0}(X)$, appears more than  
\[
\kappa' = \frac{\kappa}{\ell_1} = \varepsilon_\text{R1CS}\cdot |\mathbb F_q|^t
\]
many times in the list. 
This $p_1(X)$ has more than $\kappa'$ many $\vec\xi_i, p_{2,i}(X)$ satisfying answers for the lifted randomized R1CS on input $\vec{in}$. 
Since the soundness error of the lifted R1CS equals that of the original R1CS,  $p_1(X)$ is a solution of the first stage R1CS. 
This proves that $f_1\in \mathscr R_{\theta_1}$.
$\Box$
\end{proof}

\paragraph{Phase \ref{i:main:phase:2}:}
The second phase, the zero-check, reduces $\mathcal R_\text{phase 1}$ to 
$\mathcal R_\text{phase 2}$, where 
\begin{equation*}
\mathcal R_\text{phase 2} =
\left\{
    \begin{matrix}
        f_1 \in \left(\mathbb F_q^{D_1}\right)^{e_1}, %v_1\in F 
        \\
        f_2 \in \left(\mathbb F_q^{D_k}\right)^{e_k}, %v_2\in F 
        \\
        q_1(X), \ldots, q_m(X) \in F[X]^{\leq 3}
        \\
        v_A, v_B, v_C \in F
    \end{matrix}
    \::\:
    \hspace*{-10mm}
    \begin{matrix}
        \exists (p_1(X),p_2(X)) \in \mathscr L_{m_1} \times \mathscr L_{m_2}
        \\
        d(p_1, f_1) \leq \theta_1, d(p_2, f_2)\leq \theta_2
        % \\
        % p_1(z_1) = v_1, p_2(z_2) = v_2,
        \\\text{and}
        \\
        0 = q_1(0) + q_1(1),
        \\
        q_i(\alpha_i) = q_{i+1}(0) + q_{i+1}(1),
        \\
        \text{for }i = 1,\ldots, m-1,
        \\
        q_m(\alpha_m) = (v_A \cdot v_B - v_C)\cdot eq(\vec\alpha, \vec\rho)
    \end{matrix}
\right\},
\end{equation*}
depends on the verifier challenges $\vec\rho,\vec\alpha$ (and the preceding $\vec\xi$). 
The soundness error of this reduction is 
\[
\varepsilon_\text{phase 2}  \leq \ell_1\cdot \ell_2 \cdot
    \left(
    \frac{m}{|F|}
    + m\cdot \frac{3}{|F|}
    \right), 
\]
again scaled by the number of joint $\theta_1$- and $\theta_2$-proximates of $f_1$ and $f_2$.
The components correspond to the roundwise soundness errors of that phase. 
The first term is the probability of that $\vec \rho$ is a zero of the multilinear extension of $a(\vec x)\cdot b(\vec x) - c(\vec x)$, $\vec x \in H_m$, despite being non-trivial.   
The second term is for the $m$ rounds of the sumcheck. 
Each $\frac{3}{|F|}$ stands for the probability that the specialization of the sumcheck claim to the random $\alpha_i$ is satisfied by a joint proximate $(p_1(X), p_2(X))$,
%$\in \mathscr L_{m_1}\big|_{z_1,v_1}\times \mathscr L_{m_2}\big|_{z_1,v_1}$, 
although not both claims for $X_i = 0$ and $X_i= 1$ hold. 

\paragraph{Phase \ref{i:main:phase:3}:}
The third phase reduces $\mathscr R_{phase 2}$ to the final verifier relation $\mathscr R_V$ of the protocol. 
This is 
\[
\mathscr R_{V} =
\left\{
    \begin{matrix}
        f_1 \in \left(\mathbb F_q^{D_1}\right)^{e_1}, %v_1\in F 
        \\
        f_2 \in \left(\mathbb F_q^{D_k}\right)^{e_k}, %v_2\in F 
        \\
        q_1(X), \ldots, q_m(X) \in F[X]^{\leq 3}
        \\
        v_A, v_B, v_C \in F
        \\
        g_2 \in \left(\mathbb F_q^{D_{2}}\right)^{e_{2}}, u_2\in F 
        \\
        \vdots
        \\
        g_{r-1} \in \left(\mathbb F_q^{D_{r}}\right)^{e_{r}}, u_{r-1}\in F
        \\
        g_r \in \mathscr L_{k_r} 
    \end{matrix}
    \::\:
    % \hspace*{-5mm}
    \begin{matrix}
        % \exists (p_1(X),p_2(X)) \in \mathscr L_{m_1}\big|_{z_1,v_1}\times \mathscr L_{m_2}\big|_{z_1,v_1}
        % \\
        % d(p_1, f_1) \leq \theta_1, d(p_2, f_2)\leq \theta_2
        % % \\
        % % p_1(z_1) = v_1, p_2(z_2) = v_2,
        % \\\text{and}
        \\
        0 = q_1(0) + q_1(1),
        \\
        q_i(\alpha_i) = q_{i+1}(0) + q_{i+1}(1),
        \\
        \text{for }i = 1,\ldots, m-1,
        \\
        q_m(\alpha_m) = (v_A \cdot v_B - v_C)\cdot eq(\vec\alpha, \vec\rho)
        \\\text{and}
        \\
            g_r \text{ satisfies the aggregated WHIR claim}
        \\
        \bar v = \langle K_r, g_r\rangle_{H_{k_r}}, 
        \\
        \text{which involves the specializations}
        \\
        M_i(\vec\alpha, \vec\beta \| \: .\:), 
        \\ 
        M=A,B,C, \: i=1,2
    \end{matrix}
\right\},
\]
where $g_2$, \ldots, $g_{r-1}$, are the folding proposals committed by in the first $r-1$ rounds of WHIR, together with their values $u_2$, \ldots, $u_{r-1}$ at the verifier DEEP queries $z_2,\ldots, z_{r-1}$, and $g_r$ is the final folding proposal revealed in plain. 

The soundness error of this phase is that of joint WHIR, 
\[
\varepsilon_\text{phase 3} = \varepsilon_{WHIR}(\theta_1,\ldots, \theta_r)
\]
and will be described in the following section.



\subsection{Soundness ZO0K. \textcolor{red}{(tbc)}}

For now, only few notes:

For simplicity, we assume the masks as given. 
(The full story is a bit more complicated as it needs to take into account the proximates of the masks.)
Suppose that with notable probability in $\vec\beta_i$ the prover is able to provide $\hat f_i$ which close to a poly $\hat p_i$, so that
\begin{itemize}
    \item 
    the new ensemble with $\hat p_i$ and $\msk_i$ satisfies the $\vec\beta_i$-specialized claim \eqref{e:zook:specialized:claim},
    \item 
    the folding $\mathsf{fold}_{\vec\beta_i}\hat f_{i-1}$ agrees on a density $\alpha$ subset of $D_{i-1}$ with  $\hat p_i - \msk_i$.
    (Which is caused by \eqref{e:zook:out:domain:constraint} and \eqref{e:zook:in:domain:constraint} satisfied with notable prob.)
\end{itemize}

Then, with the help of line decodability, all intermediate foldings are proximate to polys, the mv rep of which satisfy the specialized claim, up to that $\hat f_{i-1}$ has a close poly $\hat p_{i-1}$, satisfying the combined claim \eqref{e:zook:claim:combined}.



\subsection{The complete protocol \textcolor{red}{(tbc)}}